\section*{Chapter \ref{ch:m1:volatility-first-contact} --- Volatility as a Traded Quantity: First Contact}

\begin{solution}{exo:m1:volatility-first-contact:1}
$F = 100 + \mathrm{e}^{0.01}(4.06 - 3.87) = 100.19$. Implied carry: $100 -
\mathrm{e}^{-0.01} \times 100.19 = 0.81$: an expected dividend of about eighty
cents (one strike with prices rounded to the cent is good to a cent or two).
\end{solution}

\begin{solution}{exo:m1:volatility-first-contact:2}
The strike is slightly below the forward, so the call is a little in the
money and the put a little out: use their average, 3.965. $\sigma \approx
3.965/(0.4 \times 100.2 \times 0.990 \times 0.5) = 20.0\%$, against 20.1 from
the solver.
\end{solution}

\begin{solution}{exo:m1:volatility-first-contact:3}
$2.55 \times \$1\,000 \times 40 = \$102\,000$; $2.55/0.05 = 51$ ticks.
\end{solution}

\begin{solution}{exo:m1:volatility-first-contact:4}
The call's lower bound is $\mathrm{e}^{-rT}(F - K) = 0.990 \times 10.20 =
10.10$. A price of 10.05 is below it: no volatility, however small, produces
it. The system must flag the quote (stale, crossed with the share, or the
forward is wrong, for instance a dividend not yet in the data) and exclude it
from any surface; returning zero or the previous value hides the problem.
\end{solution}

\begin{solution}{exo:m1:volatility-first-contact:5}
Carry $= 100 - \mathrm{e}^{-0.01} \times 98.10 = 2.88$ over a quarter: 11.5\% a
year. With no dividend, it is the borrow fee that the marginal arbitrageur
pays to be short: the share is special, and puts look expensive against
calls by exactly that amount. It is not an arbitrage for anyone who must pay
the same fee.
\end{solution}

\begin{solution}{exo:m1:volatility-first-contact:6}
$\sum r^2 = 0.16 + 1.21 + 0.49 + 4.00 + 2.25 = 8.11$ (in \%$^2$); mean 1.622;
daily volatility 1.274\%; annualised $\times\sqrt{252}$: 20.2\%. Over these
five days the buyer of options at 16 was right; five days are far too few to
say more.
\end{solution}

\begin{solution}{exo:m1:volatility-first-contact:7}
With cent prices: average forward 100.20, strike-by-strike range under three
cents, implied dividend 0.80, skew 6.2 points. With five-cent prices the
forward and dividend move by a few cents (the dividend estimate, 0.80, has an
error of the same size as the rounding, several percent of itself) while the
near-the-money volatilities barely change; the volatilities of the cheapest
options, worth 0.05 to 0.35, become useless, since five cents is a large
fraction of their price. Wide markets blur the wings first.
\end{solution}

\begin{solution}{exo:m1:volatility-first-contact:8}
The index cannot be bought; the future can, at 15.6, and the curve already
contains the mean reversion: the second future is at 16.7, the seventh at
19.0. If the index stays at 14, the future bought at 15.6 converges to 14 and
loses 1.6 points, \$1\,600 a contract, in a month, every month. Buying is a
bet that the index rises \emph{faster than the curve says}, against a carry
that the sellers of volatility collect the rest of the time.
\end{solution}

\begin{solution}{pb:m1:volatility-first-contact:1}
\textbf{1.} $F_K = K + \mathrm{e}^{0.01}(C - P)$: 100.20, 100.20, 100.19,
100.20, 100.20.
\textbf{2.} Prices are rounded to the cent and multiplied by 1.01. The least
reliable are the strikes where one of the two options is far in the money:
its quote is wide and it may carry an early-exercise premium if American.
\textbf{3.} 100.20.
\textbf{4.} $100 - 0.990 \times 100.20 = 0.80$.
\textbf{5.} No: the implied carry equals the expected dividend (0.80, worth
0.797 today). Nothing is left over for a borrow fee.
\textbf{6.} The out-of-the-money option: the put below the forward, the call
above. It has the tighter quote relative to its sensitivity to volatility,
no intrinsic value to subtract, and no early-exercise premium to speak of.
\textbf{7.} $23.8 - 17.6 = 6.2$ points; 3.1 points per 10\% of strike.
\textbf{8.} $(4.06 + 3.87)/2 = 3.965$; $3.965/(0.4 \times 100.2 \times 0.990
\times 0.5) = 20.0\%$.
\textbf{9.} With $F = 100$ in place of 100.20 the call's intrinsic value is
understated by 0.20, so more of its 11.22 is attributed to volatility: 25.1
instead of 23.8. In-the-money options have little sensitivity to volatility,
so a small error in the forward becomes a large error in \omterm{def:m1:volatility-first-contact:iv}{implied
volatility}; and the call and the put at the same strike then disagree, which
is the symptom to look for.
\textbf{10.} $4.03 - 3.90 - 100.01 + 0.797 + 99.005 = -0.08$.
\textbf{11.} $3.84 - 4.09 + 99.99 - 0.797 - 0.05 - 99.005 = -0.11$.
\textbf{12.} Both lose: parity holds inside the bid-ask spreads, and the
chain is consistent with one forward. That is the normal state, maintained by
the desks for which these two numbers are closest to zero.
\textbf{13.} A conversion would make $4.60 - 3.90 - 100.01 + 0.797 + 99.005 =
+0.49$ a share. First check that the quote is real and current (size,
exchange, time), that the share price is current, that no dividend or
corporate action has changed the contract, and that the option is not about
to be adjusted; then sell it.
\textbf{14.} Not by the price. With the dividend unchanged the forward is
97.17, and 2.74 is an \omterm{def:m1:volatility-first-contact:iv}{implied volatility} of 20.5 against 20.1 the day before:
in volatility terms the call is slightly dearer. One compares volatilities,
and beyond that, volatilities at comparable moneyness.
\textbf{15.} Because volatility is stable when the share moves, comparable
across strikes, expiries and underlyings, and is the quantity a hedged
option position is exposed to.
\textbf{16.} That the market pays more, per unit of model risk, for
protection against falls than for participation in rises: fear of jumps
downwards, demand for puts from holders of the share, supply of calls from
overwriters.
\textbf{17.} It weights options by $1/K^2$ and includes the whole strip: the
expensive low-strike puts count most, so the number is an average of the
smile tilted towards its high side.
\textbf{18.} Sellers of options bear a risk that pays off badly in bad times;
they are compensated by a premium, like an insurer. Implied above realised on
average is that premium, and it is paid back, with interest, in months like
February 2018.
\textbf{19.} \textbf{An implied dividend of 0.80 and a 90--110 skew of 6.2
volatility points.}
\textbf{20.} The price of an option, expressed in the one unit that stays
still when the underlying moves.
\end{solution}

\begin{solution}{iq:m1:volatility-first-contact:1}
$C - P = \mathrm{e}^{-rT}(F - K)$ for European options of the same strike and
expiry. It assumes only that one can trade the forward: buy or short the
underlying and borrow or lend cash. Dividends and the borrow fee enter
through $F$. No model, no distribution. For American options it becomes a
pair of inequalities.
\iqlookfor{model-free; the forward, not the spot; the American caveat.}
\end{solution}

\begin{solution}{iq:m1:volatility-first-contact:2}
The volatility that makes the Black--Scholes formula match the market price:
a unit conversion. It contains the expected volatility, a risk premium for
bearing it, and supply and demand for that strike; it differs by strike,
which no single forecast could. On average it has exceeded subsequent
\omterm{def:m1:volatility-first-contact:realised}{realised volatility}.
\iqlookfor{the risk premium, and the smile as proof that it is not a
forecast.}
\end{solution}

\begin{solution}{iq:m1:volatility-first-contact:3}
Check the premium against its bounds (discounted intrinsic and the
discounted forward or strike) and refuse outside them. Use the
out-of-the-money option via parity. Bisection always converges; Newton is
faster but fails where vega is tiny, deep in or out of the money, so
safeguard it with brackets. Beware: wrong forward (dividends), American
premium, expiry-day options with $T \to 0$, prices of one tick where the
answer spans twenty volatility points, and time conventions (calendar or
business time).
\iqlookfor{bounds check, vega collapse, and the forward as the main source
of error.}
\end{solution}

\begin{solution}{iq:m1:volatility-first-contact:4}
A dividend I have not included, or a larger one than I assumed. A borrow fee:
hard-to-borrow shares have low forwards. The interest rate I use. \omterm{def:m1:option-contracts-and-exchanges:style}{American
exercise}: the put carries an early-exercise premium. Stale or wide quotes
compared at mid. An expected corporate action changing the deliverable.
Only when the conversion or reversal is profitable at the prices I can
actually trade, including borrow, is it an arbitrage.
\iqlookfor{dividends and borrow before anything else.}
\end{solution}

\begin{solution}{iq:m1:volatility-first-contact:5}
From out-of-the-money index options of two expiries bracketing thirty days:
each expiry's \omterm{def:m1:volatility-first-contact:realised}{variance} is a sum of mid-quotes weighted by $\Delta K/K^2$, with
a small correction for the distance between the forward and the nearest
strike; the two \omterm{def:m1:volatility-first-contact:realised}{variances} are interpolated to thirty days, and the index is a
hundred times the square root. It is the price of a portfolio that would
have to be rebalanced continuously and whose composition changes every
moment; one can trade futures and options on it, which price its future
value, not the index itself.
\iqlookfor{the strip and its weights; futures are not the index.}
\end{solution}

\begin{solution}{iq:m1:volatility-first-contact:6}
Products short the front volatility futures had earned the roll for years and
grown large. They rebalance daily to constant leverage: after a rise in the
futures, both inverse and leveraged-long products must buy, in the last
minutes, an amount proportional to the move. On 5 February a 4\% fall in the
index raised the futures; the required purchases were large relative to the
futures' liquidity, pushed the futures up further, and an inverse product
lost nearly everything in a session. Visible beforehand: the products' size
relative to futures volume, the convexity of the rebalancing rule (demand
grows with the move), the concentration of that demand at one time of day,
and a position whose worst day exceeds its total historical profit.
\iqlookfor{the rebalancing arithmetic and size relative to liquidity, not
just ``volatility spiked''.}
\end{solution}
